\chapter*{\sffamily Introducción}

\addcontentsline{toc}{chapter}{\sffamily Introducción}
\markboth{\sffamily Introducción}{}

\par Cada semestre, cerca de mil estudiantes de la sede Palmira de la Universidad
Nacional de Colombia cursan Cálculo Diferencial o Matemáticas Básicas. Para la
mayoría es el comienzo de la carrera. Para una parte, es también el primer
tropiezo serio: son asignaturas con pérdidas históricamente altas en la sede, y
para un estudiante que empieza, perderlas una y otra vez puede terminar en la
salida de la universidad. A esas mismas aulas llegó, en los últimos años, la
inteligencia artificial generativa. Llegó sin diseño pedagógico: los estudiantes
ya le piden a Gemini o a ChatGPT que les resuelva los ejercicios, con o sin
autorización de sus docentes.

Este trabajo nace en un lugar preciso: la tutoría personalizada del programa de
Grupos de Estudio Autónomo (GEA) de la sede, adonde llegan los estudiantes de
esas dos asignaturas que sus docentes, otros tutores o el área de acompañamiento
académico y psicosocial remiten por sus dificultades. Como tutor de ese
componente, el autor ve de cerca un patrón que la investigación sobre las
matemáticas universitarias lleva años describiendo: estudiantes que avanzan
cuando alguien les muestra el procedimiento y que todavía no logran resolver
solos. Una IA que entrega el procedimiento completo puede aliviar esa dificultad
o agravarla. ¿De qué depende? La hipótesis de trabajo es que depende, sobre todo,
de la tarea en la que se inserta.

Por eso el objeto de estudio es el diseño de tareas. La propuesta consiste en
diseñar una secuencia de tareas de resolución de problemas de cálculo diferencial,
centrada en el límite, en la que cada problema se resuelve primero sin IA, luego
con Gemini en un rol de acompañamiento y al final con una validación en la que el
estudiante compara ambos momentos. La implementación sigue la lógica del
experimento de enseñanza, en las sesiones de tutoría personalizada, y el análisis
es cualitativo: se busca comprender qué estrategias, razonamientos y errores
aparecen con y sin IA, y qué tipo de apoyo encuentran los estudiantes en ella.

Este documento presenta un avance de la propuesta de trabajo final de maestría,
en la modalidad de profundización, de la Maestría en Enseñanza de las Ciencias
Exactas y Naturales de la Universidad Nacional de Colombia, sede Palmira. Su
objetivo general es diseñar, implementar y analizar una secuencia de tareas
mediada por inteligencia artificial generativa, orientada a la resolución de
problemas de cálculo diferencial en estudiantes de Cálculo Diferencial y
Matemáticas Básicas remitidos a tutoría personalizada.

El \textbf{Capítulo 1} plantea el problema a partir del contexto de la sede, del
programa GEA y de los registros de la tutoría personalizada; formula la pregunta
de investigación, sus preguntas derivadas y los objetivos. El \textbf{Capítulo 2}
reúne el marco teórico y el marco contextual, que describe la sede, el programa y
las condiciones tecnológicas en que se usa la IA. El \textbf{Capítulo 3} expone la
metodología: el experimento de enseñanza, los participantes, la secuencia de
tareas, el papel de Gemini, los instrumentos ---entre ellos una entrevista
semiestructurada basada en tareas---, el análisis por objetivo específico, el
cronograma, el presupuesto y las consideraciones éticas. El \textbf{Anexo A}
contiene el manual de uso de Gemini para los estudiantes, producto del segundo
objetivo específico.

\clearpage
\newpage
